
%%%%%%%%%%%%%%%%%%%%%%%%%%
\documentclass[12pt, a4paper]{article}
\input{systeme}
\usepackage[a4paper, top=4cm,left=2.5cm, right=2.3cm, bottom=2.7cm]{geometry}
\usepackage{amsmath, amssymb, array, multicol, enumitem}
\usepackage{pdfpages}
\usepackage{comment, graphicx}  
\usepackage{xparse}
\DeclareGraphicsExtensions{.png,.jpg}
\usepackage{pgffor}  % for loops in latex
\usepackage{ifthen}  % conditional latex compilation
\usepackage{todonotes}  % notes in bright orange
\usepackage[colorlinks=true]{hyperref}  % hyperlinks
\AtBeginDocument{% coloring of links
	\hypersetup{pdfborder={0 0 0},
		urlcolor=blue}%
}
\usepackage{fontspec}  % Hebrew fonts
\usepackage{xunicode}  % Hebrew character encoding
\usepackage{polyglossia}  % Hebrew Latex
\usepackage[absolute]{textpos}   % Proper positioning of images, e.g. BGU logo

\setdefaultlanguage{hebrew}
\setotherlanguage{english}

\newfontfamily\hebrewfont[Script=Hebrew]{Times New Roman}

\makeatletter
\def\tagform@#1{\maketag@@@{(\ignorespaces#1\unskip)}}
\makeatother


\makeatletter
\renewcommand{\SepMark}[1]{\def\@SepMark{#1}}
\makeatother
\SepMark{.}
\robustify\hebrewnumeral
\robustify\Hebrewnumeral
\robustify\Hebrewnumeralfinal

\def\urladdrname{{\itshape כתובת רשת}}
\def\emailaddrname{{\itshape כתובת דואל}}

\makeatletter
\let\orig@MakeUppercase \MakeUppercase
\makeatother
\providecommand{\MakeUppercase}[1]{#1}

\fontsize{14}{15}  

\setdefaultlanguage{hebrew}
\setotherlanguage{english}
\thispagestyle{empty}
\selectlanguage{hebrew}

\newcommand{\cc}{\mathbb{C}}
\newcommand{\CC}{\mathbb{C}}
\newcommand{\rr}{\mathbb{R}}
\newcommand{\RR}{\mathbb{R}}
\newcommand{\nn}{\mathbb{N}}
\newcommand{\zz}{\mathbb{Z}}
\newcommand{\ZZ}{\mathbb{Z}}
\newcommand{\qq}{\mathbb{Q}}
\newcommand{\FF}{\mathbb{F}}
\newcommand{\QQ}{\mathbb{Q}}
\DeclareMathOperator{\vol}{vol}
\newcommand{\abs}[1]{\left|#1\right|}

\DeclareMathOperator{\Adj}{Adj}
\DeclareMathOperator{\deter}{det}
\DeclareMathOperator{\Img}{Im}
\DeclareMathOperator{\Ker}{Ker}
\DeclareMathOperator{\rank}{rank}
\DeclareMathOperator{\Span}{Span}
\DeclareMathOperator{\tr}{tr}


\newcommand{\calB}{\mathcal{B}}
\newcommand{\calC}{\mathcal{C}}
\newcommand{\cD}{\mathcal{D}}
\newcommand{\calE}{\mathcal{E}}


%%%% new for complex functions

\newcommand{\ol}{\overline}
\newcommand{\wt}{\widetilde}
\newcommand{\wh}{\widehat}
\newcommand{\re}{\mathrm{Re}}
\newcommand{\im}{\mathrm{Im}}
\newcommand{\goes}{\longrightarrow}
\newcommand{\cO}{\mathcal{O}}
\newcommand{\cR}{\mathcal{R}}
\newcommand{\Log}{\mathrm{Log}}
\newcommand{\PP}{\mathbb{P}}

%%%%%probabilty shorts

\newcommand{\Geo}{\mathrm{Geo}}
\newcommand{\Bin}{\mathrm{Bin}}
\newcommand{\Poi}{\mathrm{Poi}}
\newcommand{\Ber}{\mathrm{Ber}}
\newcommand{\Var}{\mathrm{Var}}
\newcommand{\EE}{\mathbb{E}}




%%%%%%%%%%%%%%%%%%%%%%%%%%%%%%%%%%%%


\begin{document}
\selectlanguage{hebrew}

\newcommand{\semester}{%
	סמסטר א' תשפ״ה
}
\newcommand{\course}{%
הסתברות לסטטיסטיקה}
\newcommand{\coursenumber}{9321}
\newcommand{\doctype}{%
\boxed{\text{
 דף  תרגילים 3:
פתרון}}
}
\newcommand{\hwknumber}{}
\newcommand{\chaptype}{%
	פרק }
\newcommand{\chapnum}{\textenglish{I}}

\newcommand{\chapname}{%
אאא } 
	
\thispagestyle{empty}
\begin{textblock*}{9cm}[0.5,0.5](90mm,3cm)
	\includegraphics[width=2in]{logo}
\end{textblock*}	

\begin{center}
	\large	
המחלקה למתמטיקה - 
	\semester  \\
	{\large \bf \course - 201.1.9921}\\
	\medskip
	{\Large  \textbf{\doctype \hwknumber}  }\\
\end{center}
\vspace{0.3cm}

%%%%%%%%%%%%%%%%%%%%%%%%%%%%%%%%%%%%

\begin{enumerate}

\item
שנה אחרי שנה פרופסור אקס מלמד את אותו הקורס ״מבוא למוטנטים״ באוניברסיטה.  
הקורס נחשב לקשה וידוע כי ההסתברות שיש סטודנט (אחד לפחות) שקיבל
$100$
בקורס היא
$0.02$,
כך זה בכל שנה באופן בלתי תלוי בשנים הקודמות בהן התקיים הקורס.
פתרו כל אחד מהסעיפים הלא קשורים הבאים.
\begin{enumerate}
\item
מה ההסתברות שרק בשנה הרביעית היה סטודנט שקיבל
$100$,
לראשונה בתולדות הקורס?
\item
מה ההסתברות שב-$6$
השנים הראשונות לא היה סטודנט שקיבל
$100$
כלל?
\item
מה ההסתברות שב-$10$ השנים הראשונות, היו בדיוק שנתיים בהן היה סטודנט שקיבל 
$100$?
\item
 בכל שנה,
 פרופ׳ מקגונגל (קולגה של פרופ׳ אקס) מלמדת במקביל אליו את הקורס ״מבוא לקסמים״.
 
ידוע כי  ההסתברות שיש לפחות סטודנט אחד שקיבל בקורס שלה
$100$
היא
$0.05$,
באופן בלתי תלוי בשנים הקודמות בהן התקיים הקורס.

עוד ידוע כי הקורסים של פרופ׳ אקס ופרופ׳ מקגונגל  בלתי תלויים אחד בשני.

שני הפרופסורים הסכימו (כל אחד/ת עבור הקורס שלו/ה) כי בשנה הראשונה  בה יש סטודנט המקבל
$100$
בקורס, הוא/היא יוצא/ת לשנת שבתון. 
\begin{enumerate}
\item[(i)]
מה ההסתברות כי שניהם יצאו לשבתון לאחר שנה בלבד?
\item[(ii)]
מה ההסתברות כי שניהם יצאו לשבתון בדיוק באותה השנה? 
\end{enumerate}
\end{enumerate}

\noindent\rule{15cm}{0.4pt}

\underline{פתרון:}
נסמן ב-$X$
את מספר השנים שלקח עד שהיה סטודנט שקיבל
$100$
בקורס. לכן
$X\sim \Geo(0.02)$.
\begin{enumerate}
\item
בסעיף זה שואלים בעצם לגבי ההסתברות הבאה
$$\PP(X=4)
(1-0.02)^3\cdot 0.02=0.98^3\cdot 0.02=
0.0188$$
%%%%

\noindent\rule{15cm}{0.4pt}

\item
בסעיף זה שואלים על המאורע בו ב-$6$
השנים הראשונות לא היה סטודנט שקיבל 
$100$
(כלומר, לא הייתה הצלחה $6$
שנים), כלומר
$$\PP(X>6)=(1-0.02)^6=0.98^6=0.8858$$
%%%%

\noindent\rule{15cm}{0.4pt}

\item
כאן שואלים על מספר ההצלחות אבל בטווח של
$10$
שנים, זוהי כבר התנהגות של משתנה בינומי, אז נסמן ב-$Y$
את מספר השנים בהן היה סטודנט שקיבל
$100$
בקורס, מתוך 
$10$
השנים הראשונות. לכן
$Y\sim \Bin(10,0.02)$
ונחשב
$$\PP(Y=2)
=\binom{10}{2}\cdot
(0.02)^2\cdot(0.98)^8=0.0153$$
%%%%
\item
ניעזר ב-$X$
אותו הגדרנו בתחילת השאלה, ונגדיר משתנה נוסף
$Z$
להיות מספר  השנים שלקח עד שסטודנט קיבל
$100$
בקורס של פרופ׳ מקגונגל. אז מהנתונים נסיק כי
$Z\sim\Geo(0.05)$
וכי
$X,Z$
משתנים בלתי תלויים.
\begin{enumerate}
\item[(i)]
כאן  השאלה היא מהי ההסתברות של המאורע שהוא
"$X=1$
וגם
$Z=1$":
$$\PP(X=Z=1)=\PP(X=1)\cdot \PP(Z=1)$$
כי הם בלתי תלויים, ונמשיך לחשב
$$\PP(X=Z=1)
=0.02\cdot 0.05
=0.001$$
\item[(ii)]
נסמן ב-$A$
את המאורע המבוקש ״שניהם יצאו לשבתון באותה השנה״, ונשים לב כי
$$A=
\{X=Z=1\}\cup\{X=Z=2\}\cup\ldots=
\bigcup_{n=1}^{\infty}
A_n$$
כאשר
$A_n=\{ X=Z=n\}$
זה בעצם המאורע ששניהם יצאו לשבתון בדיוק בשנה ה-$n$.
בגלל שהמאורעות האלו
$A_1,A_2,\ldots$
זרים בזוגות, אז
\begin{align*}
\PP(A)=\PP(A_1\cup A_2\cup\ldots)=
\sum_{n=1}^{\infty}
\PP(A_n)
=\sum_{n=1}^{\infty}=
\end{align*}
לכל
$n\ge1$
נשתמש בעובדה כי
$X,Y$
משתנים בלתי תלויים ואז
$$\PP(A_n)=\PP(X=n)\cdot \PP(Y=n)
=0.98^{n-1}\cdot 0.02\cdot 0.95^{n-1}\cdot 0.05
=q^{n-1}\cdot 0.001$$
כאשר
$q=0.98\cdot 0.95$.
לכן
$$\PP(A)=0.001\cdot
\sum_{n=1}^{\infty}
q^{n-1}
=0.001\cdot
\frac{1}{1-q}=0.0144$$
כאשר נעזרנו בסכום של סדרה הנדסית אינסופית.
\end{enumerate}
\end{enumerate}

\newpage
\item
יהיו
$X\sim \Bin(4,0.3)$
ו-
$Y\sim \mathrm{Poi}(3)$
שני משתנים מקריים
ונניח כי הם
 בלתי תלויים.

חשבו את ההסתברויות הבאות:
\begin{enumerate}
\begin{multicols}{3}
\item
$\PP(X+Y=1)$.
\item
$\PP(X+Y=5)$.
\item
$\PP(X< Y)$.
\end{multicols}
\end{enumerate}

\noindent\rule{15cm}{0.4pt}

\underline{פתרון:}
נזכיר כי טווח הערכים של 
$X$
 הוא
 $\cR_X=
 \{0,1,2,3,4\}$
 וטווח הערכים של
 $Y$
 הוא
 
 $\cR_Y=\{0,1,2,\ldots\}$.
\begin{enumerate}
\item
נקבל
 \begin{align*}
 \PP(X+Y=1)
 &=
 \PP(X=0, Y=1)
 +\PP(X=1,Y=0)
 \\&=\PP(X=0)\cdot \PP(Y=1)
 +\PP(X=1)\cdot \PP(Y=0)
 \\&=
 0.7^4\cdot e^{-3}\cdot \frac{3^1}{1!}+
 \binom{4}{1}0.3\cdot 0.7^3\cdot e^{-3}\cdot
 \frac{3^0}{0!}
 \end{align*}
 %%%
 
 \noindent\rule{15cm}{0.4pt}

 \item
 באופן דומה לסעיף הראשון, נציג את המאורע
 $\{X+Y=5\}$
 כאיחוד של מספר מאורעות זרים בזוגות ונשים לב כי
 $X=4$
 הערך הכי גדול של
 $X$
 כי המאורע
 $X=5$
 הוא ריק, אז נקבל
 $$\PP(X+Y=5)=
\sum_{k=0}^5
 \PP(X=k, Y=5-k)
 =\sum_{k=0}^4
 \binom{4}{k}0.3^k0.7^{4-k}\cdot
 e^{-3}\cdot
 \frac{3^{5-k}}{(5-k)!}$$
 %%%%
 
 \noindent\rule{15cm}{0.4pt}

 \item
 גם כאן נחשוב על המאורע המבוקש
 $\{X<Y\}$
 ונבין כי הוא מורכב מהרבה מאוד (אינסוף) מקרים מהצורה
 $\{X=k, Y=m\}$
 כאשר
 $k=0,1,2,3,4$
 וגם
 $n>k$.
 במקום לחשב זאת באופן ישיר, נעדיף לחשב את ההסתברות של המאורע המשלים
 \begin{align*}
 \PP(X<Y)
& =1-\PP(X\ge Y)
 =1-\sum_{k=0}^4
 \PP(X=k, Y\le k)
 \\&=
 1-\PP(X=0,Y=0)-\PP(X=1,Y=0)-\PP(X=1,Y=1)
\\& -\PP(X=2,Y=0)-\PP(X=2,Y=1)-\PP(X=2,Y=2)
\\& -\PP(X=3,Y=0)-\ldots-\PP(X=3,Y=3)
\\& -\PP(X=4,Y=0)-\ldots-\PP(X=4,Y=4) 
 \end{align*}
 מפרידים כל הסתברות לכפל הסתברויות כי המשתנים בלתי תלויים וממשיכים את החישוב הסיזיפי.
\end{enumerate}
\newpage

\item
זורקים
$2$
מטבעות הוגנים ובלתי-תלויים עם 
$0$
ו-$1$
משני צידי המטבעות.
נגדיר את המשתנים המקריים הבאים:
$X_0$
מספר ה-$0$-ים שיצאו.
\quad
$X_1$
מספר ה-$1$-ים שיצאו.
\quad
$Y$
סכום המספרים שיצאו במטבעות.

$Z$
שווה
$1$
אם תוצאות המטבעות שונות,
ושווה
$0$
אם תוצאות המטבעות זהות.

בכל אחד מהסעיפים הבאים בדקו
האם 
שני המשתנים הנתונים
הם בלתי תלויים או לא, בעזרת ההגדרה.
\begin{enumerate}
\begin{multicols}{4}
\item
$X_0,X_1$.
%%%
\item
$X_0, Y$.
%%%%
\item
$X_1,Z$.
%%%%
\item
$Y,Z$.
\end{multicols}
\end{enumerate}

\noindent\rule{15cm}{0.4pt}

\underline{פתרון:}

\begin{enumerate}
\item
המשתנים תלויים. נשים לב כי
$\PP(X_0=2,X_1=2)=0$
כי זהו מאורע ריק (שבשני מטבעות יצא
$0$
וגם בשני מטבעות יצא
$1$)
 אבל מצד שני
$$.\PP(X_0=2)\cdot\PP(X_1=2)=\frac{1}{4}\cdot\frac{1}{4}\neq0
=\PP(X_0=2,X_1=2)$$
%%%

\noindent\rule{15cm}{0.4pt}

\item
המשתנים תלויים. נשים לב כי
$\PP(X_0=0,Y=0)=0$
כי זהו מאורע ריק (שבשני מטבעות יצא
$1$
וגם שזכום המטבעות הוא
$0$)
 אבל מצד שני
$$.\PP(X_0=0)\cdot\PP(Y=0)
=\frac{1}{4}\cdot\frac{1}{4}\neq0
=\PP(X_0=0,Y=0)$$
%%%

\noindent\rule{15cm}{0.4pt}

\item
המשתנים תלויים. נשים לב כי
$\PP(X_1=2,Z=1)=0$
כי זהו מאורע ריק (שבשני מטבעות יצא
$1$
וגם בשני מטבעות יצאו תוצאות שונות) אבל מצד שני
$$.\PP(X_1=2)\cdot\PP(Z=1)=\frac{1}{4}\cdot\frac{1}{2}\neq0
=\PP(X_1=2,Z=1)$$

\noindent\rule{15cm}{0.4pt}

\item
המשתנים תלויים. נשים לב כי
$\PP(Y=0,Z=1)=0$
כי זהו מאורע ריק (שבשני מטבעות יצא
$0$
וגם בשני מטבעות יצאו תוצאות שונות) אבל מצד שני
$$.\PP(Y=0)\cdot\PP(Z=1)=\frac{1}{4}\cdot\frac{1}{2}\neq0
=\PP(Y=0,Z=1)$$
\end{enumerate}

\newpage

\item
נדב קיבל מתנה לכבוד יום ההולדת שלו: שני כרטיסי גירוד.

בכל כרטיס גירוד ניתן לזכות ב-$100$
ש״ח בהסתברות
$0.15$,
ב-$1000$
ש״ח בהסתברות
$0.05$
או לא לזכות כלל.

נגדיר את המשתנים המקריים:

$X$
הוא מספר הכרטיסים בהם נדב זכה בסכום כספי כלשהו (חיובי).

$Y$
הוא סכום הזכייה של נדב בסך הכל משני הכרטיסים.

\begin{enumerate}
\item
מצאו את פונקציית ההסתברות של
$X$,
את התוחלת
$\EE[X]$
ואת השונות
$\Var[X]$.
%%%
\item
מצאו את פונקציית ההסתברות של
$Y$,
את התוחלת
$\EE[Y]$
ואת השונות
$\Var[Y]$.
\end{enumerate}

\noindent\rule{15cm}{0.4pt}

\underline{פתרון:}
\begin{enumerate}
\item
מנתוני השאלה נסיק כי ההסתברות לזכות בסכום כספי כלשהו (חיובי) היא
$p=0.15+0.05=0.2$,
ולכן
$X\sim\Bin(2,p)$,
כי יש שני ניסויים בלתי תלויים (שני כרטיסים) שבכל אחד מהם ההסתברות להצלחה (שיהיה סכום כספי חיובי) היא 
$p$
והמשתנה
$X$
סופר את מספר ההצלחות (מספר הכרטיסים בהם יש סכום כספי חיובי). מהנוסחאות שלמדנו
$$.\EE[X]=2\cdot p=0.4\quad,\quad,
\Var[X]=2\cdot p\cdot (1-p)=0.32$$
%%%%

\noindent\rule{15cm}{0.4pt}

\item
נגדיר
$Y_1$
להיות סכום הזכייה של נדב בכרטיס הראשון,
ו-$Y_2$
סכום הזכייה של נדב בכרטיס השני,
לכן
$Y=Y_1+Y_2$
ונזכור כי שני המשתנים בלתי תלויים.
נכתוב את פונקציית ההסתברות של
$Y_1$
ושל
$Y_2$
שהן זהות:
$$\PP(Y_1=0)=0.8,\quad
\PP(Y_1=100)=0.15,\quad
\PP(Y_1=1000)=0.05$$
טווח הערכים של
$Y$
הוא
$$\cR_Y=
\{
0, 100, 200, 1000, 1100, 2000
\}$$
ונחשב כל אחת מההסתברויות בצורה ישירה
\begin{align*}
&\PP(Y=0)=
\PP(Y_1=0,Y_2=0)=\PP(Y_1=0)\PP(Y_2=0)
=0.8^2=0.64
\\&\PP(Y=100)=
\PP(Y_1=100)\PP(Y_2=0)+\PP(Y_1=0)\PP(Y_2=100)=2\cdot 0.8\cdot 0.15
=0.24
\\&\PP(Y=200)=\PP(Y_1=100)\PP(Y_2=100)=0.15^2=0.0225
\\&\PP(Y=1000)=
\PP(Y_1=1000)\PP(Y_2=0)+\PP(Y_1=0)\PP(Y_2=1000)=2\cdot 0.8\cdot 0.05
=0.08
\\&\PP(Y=1100)=\PP(Y_1=1000)\PP(Y_2=100)+\PP(Y_1=100)\PP(Y_2=1000)
=2\cdot 0.15\cdot 0.05=0.015
\\&\PP(Y=2000)=\PP(Y_1=1000)\PP)Y_2=1000)=0.05^2=0.0025
\end{align*}
מצאנו את פונקציית ההסתברות של
$Y$.
ניתן לחשב את התוחלת ואת השונות בצורה ישירה
$$\EE[Y]
=0\cdot 0.64
+100\cdot 0.24+
200\cdot
0.0225
+1000\cdot
0.08+
1100\cdot
0.015
+2000\cdot
0.0025
=130$$
וככה גם את התוחלת 
$$\EE[Y^2]
=100^2\cdot 0.24+
200^2\cdot
0.0225
+1000^2\cdot
0.08+
1100^2\cdot
0.015
+2000^2\cdot
0.0025=111,450$$
לכן
$$.\Var[Y]=111,450-130^2=94,550$$
את התוחלת והשונות ניתן לחשב גם בדרך אחרת:
ניתן להשתמש בלינאריות התוחלת ואז
$$\EE[Y]=\EE[Y_1+Y_2]
\EE[Y_1]+\EE[Y_2]=
2\cdot
(100\cdot 0.15+1000\cdot 0.05))=130$$
ולסיום, בגלל שהמשתנים
$Y_1,Y_2$
בלתי תלויים אפשר גם להשתמש בלינאריות השונות עבורם ולקבל
$$\Var[Y]
=
\Var[Y_1+Y_2]=\Var[Y_1]+\Var[Y_2]$$
כאשר השונות של
$Y_1$
ושל
$Y_2$
שוות, וניתנות לחישוב בצורה ישירה.

\end{enumerate}

\newpage
\item
בכל אחד מהסעיפים הלא קשורים הבאים, מצאו את פונקציית ההסתברות של
$X$,
 כאשר:
 \begin{enumerate}
%%%%%
\item
$X$
משתנה מקרי עם טווח ערכים
$\cR_X=\{0,1\}$
ונתון כי
$\EE[X]=0.25$.
%%%
\item
$X$
משתנה מקרי עם טווח ערכים
$\cR_X=\{0,1,2\}$,
ונתון כי
$\EE[X]=1.4$
וגם
$\Var[X]=0.34$.
%%%%

\end{enumerate}
\noindent\rule{15cm}{0.4pt}

\underline{פתרון:}
\begin{enumerate}
\item
מהגדרת התוחלת נקבל
$$0.25=\EE[X]
0\cdot \PP(X=0)+1\cdot \PP(X=1)=\PP(X=1)$$
ובגלל שבטווח של
$X$
יש רק שני ערכים נקבל
$\PP(X=0)=1-\PP(X=1)=0.75$.
לסיכום, קיבלנו כי
$X\sim\Ber(0.75)$.
%%%

\noindent\rule{15cm}{0.4pt}


\item
למען הנוחות נסמן
$\PP(X=0)=p_0, \PP(X=1)=p_1$
ואז
$\PP(X=2)=1-p_0-p_1$.
שוב נשתמש בהגדרת התוחלת:
$$1.4=\EE[X]
=0\cdot p_0+1\cdot p_1+2\cdot (1-p_0-p_1)=2-2p_0-p_1$$
לכן
$p_1=0.6-2p_0$.
מהגדרת השונות
$\Var[X]=\EE[X^2]-(\EE[X])^2$
 נסיק כי
$$2.3=0.34+1.96=
\EE[X^2]=0^2\cdot p_0+1^2\cdot p_1
+2^2\cdot (1-p_0-p_1)$$
נציב את
$p_1=0.6-2p_0$
ונקבל
$$2.3=0.6-2p_0+4(1-p_0-0.6+2p_0)
\Longrightarrow
p_0=0.05$$
ואז
$p_1=0.6-0.1=0.5$ וגם
$p_2=1-0.05-0.5=0.45$.
\end{enumerate}
\newpage
\item
במחלקה לאומנות יפנית עתיקה יש
$20$
סטודנטים. כל סטודנט צריך לבחור ״קורס כללי״ אחד מתוך היצע של
$100$
קורסים כלליים שהאוניברסיטה מציעה בסמסטר.
לסטודנטים לא באמת אכפת באיזה קורס כללי יבחרו, ולכן כל סטודנט בוחר באקראי אחד מבין
$100$
הקורסים.
\begin{enumerate}
\item
מהי תוחלת מספר הקורסים הכלליים בהם לא יהיה שום סטודנט מהמחלקה לאומנות יפנית עתיקה?
\item
מהי תוחלת מספר הקורסים הכלליים בהם יהיה לפחות סטודנט אחד מהמחלקה לאומנות יפנית עתיקה?
\end{enumerate}

\noindent\rule{15cm}{0.4pt}

\underline{פתרון:}
\begin{enumerate}
\item
נסמן ב-$X$
את מספר הקורסים הכלליים בהם לא יהיה שום סטודנט מהמחלקה לאומנות יפנית עתיקה, ונרצה לחשב את התוחלת של
$X$.
בשביל זה נגדיר כעת
$100$
משתנים מקריים חדשים כך:
$$X_k=
\begin{cases}
1&\text{אם בקורס ה-$k$ אין סטודנט מהמחלקה לאומנות יפנית עתיקה}
\\
0&\text{בכל מקרה אחר}\\
\end{cases} $$
לכל
$k=1,2,\ldots,100$.
נשים לב כי טווח הערכים שלו הוא
$\{0,1\}$,
לכן
$X_k\sim \Ber(p_k)$,
כאשר
$$p_k=\PP(\text{אם בקורס ה-$k$ אין סטודנט מהמחלקה לאומנות יפנית עתיקה})=
\left(\frac{99}{100}\right)^{20}=0.818$$
כי זה בעצם אומר שכל סטודנט, מבין ה-$20$, יכול לבחור כל קורס אחר פרט לקורס ה-$k$.
מהגדרת המשתנים המקריים, מתקיים
$$X=X_1+\ldots+X_{100}$$
ולכן מלינאריות התוחלת
$$.\EE[X]=\EE[X_1]+\ldots+\EE[X_{100}]=100\cdot 0.818=81.8$$
%%%%

\noindent\rule{15cm}{0.4pt}

\item
נשים לב כי 
מספר הקורסים הכלליים בהם יהיה לפחות סטודנט אחד מהמחלקה לאומנות יפנית עתיקה שווה
$100-X$,
ולכן מלינאריות התוחלת ומסעיף א׳:
$$.\EE[100-X]
=100-\EE[X]=18.2$$
\end{enumerate}
\newpage
\item
עבור
$X\sim \Bin(n,p)$
חשבו את התוחלת
$\EE[X\cdot(X-1)]$
 בשתי דרכים שונות:
 \begin{enumerate}
 \item
 (בדרך הארוכה) בעזרת כלל לוטוס וכתיבה פורמלית של התוחלת, וזה מתחיל ככה:
 $$\EE[X(X-1)]
 =\sum_{k=0}^n
 k(k-1)
 \cdot
 \PP(X=k)$$
 ניתן להשתמש כאן בזהות בה השתמשנו בכיתה
 $p\cdot
 \binom{q}{p}
 =q\cdot \binom{q-1}{p-1}$.
 \item
 (בדרך הקצרה) בעזרת לינאריות התוחלת ושימוש בתוחלת ובשונות של
 $X$.
 \end{enumerate}

\noindent\rule{15cm}{0.4pt}

\underline{פתרון:}
\begin{enumerate}
\item
ניעזר בכלל לוטוס ונקבל
\begin{align*}
\EE[X(X-1)]
 &=\sum_{k=0}^n
 k(k-1)
 \cdot
 \PP(X=k)
 \overset{(1)}{=}\sum_{k=2}^n
 k(k-1)
 \binom{n}{k}
 p^k(1-p)^{n-k}
 \\&\overset{(2)}{=}
 \sum_{k=2}^n
 (k-1)n
 \binom{n-1}{k-1}
 p^k(1-p)^{n-k}
 \\&\overset{(3)}{=}n\sum_{k=2}^n 
(n-1) \binom{n-2}{k-2}
p^k(1-p)^{n-k}
\\&= n(n-1)p^2\sum_{k=2}^n
\binom{n-2}{k-2}
p^{k-2}(1-p)^{n-2-(k-2)}
\\&\overset{(4)}{=}
n(n-1)p^2\sum_{\ell=0}^{n-2}
\binom{n-2}{\ell}p^{\ell}(1-p)^{n-2-\ell}
\overset{(5)}{=}n(n-1)p^2
\end{align*}
הסברים למעברים:

$(1)$:
עבור
$k=0,1$
המחוברים מתאפסים ולכן הסכום מתחיל מ-$2$.

$(2)$:
השתמשנו בזהות
$k\binom{n}{k}=
(n-1)\binom{n-1}{k-1}$.

$(3)$:
השתמשנו באותה הזהות כאשר
$(k-1)\binom{n-1}{k-1}=
(n-1)\binom{n-2}{k-2}$.

$(4)$:
עשינו שינוי אינדקס
$\ell=k-2$.

$(5)$:
הסכום שווה
$1$
כי זהו סכום כל ההסתברויות של משתמנה מסוג
$\Bin(n-2,p)$.

\noindent\rule{15cm}{0.4pt}

\item
בעזרת הלינאריות של התוחלת נקבל
$$\EE[X\cdot (X-1)]=
\EE[X^2-X]
=\EE[X^2]-\EE[X]=
\Var[X]+(\EE[X])^2-\EE[X]$$
כעת ניזכר כי
$\EE[X]=np,
\Var[X]=np(1-p)$
אז
$$\EE[X\cdot (X-1)]
=np(1-p)+(np)^2-np
=np(1-p+np-1)=n(n-1)p^2$$
\end{enumerate}

\end{enumerate}


\end{document}





